\section{Related Work}
\label{sec:related}

\para{Linear representations in language models.}
The Linear Representation Hypothesis (\lrh) holds that high-level concepts are encoded as directions in transformer representation spaces. \citet{park2024linear} formalize three notions of linearity---subspace, measurement, and intervention---and unify them via a causal inner product, demonstrating linear structure for 26 of 27 tested concepts in GPT-J. \citet{marks2024geometry} show that truth and falsehood are linearly represented in LLaMA-13B, and that mass-mean probes (the difference-in-means direction) are more causally meaningful than logistic regression probes, with normalized indirect effects of 0.85--0.97 versus 0.33--0.52. We adopt the mass-mean probing methodology of \citet{marks2024geometry} and apply it to a novel concept: Pok\'{e}mon franchise membership. Unlike truth or language identity, this category is culturally constructed and has inherently fuzzy boundaries, making it a more challenging test of the \lrh.

\para{Representation engineering.}
\citet{zou2023representation} introduce Representation Engineering (RepE), a top-down approach that uses contrastive stimuli to extract concept directions and then reads or steers model behavior along these directions. \citet{turner2023activation} propose Activation Addition (ActAdd), which steers model behavior by adding contrastive activation vectors to the residual stream during inference. Both approaches demonstrate that linear directions can causally influence model outputs. Our work focuses on the probing (reading) side---we extract and validate a Pok\'{e}mon direction but leave causal steering experiments to future work.

\para{Factual knowledge in transformers.}
\citet{meng2022locating} show that factual associations in GPT-2 are stored in middle-layer MLP modules and can be edited with rank-one updates. Their causal tracing technique reveals that the last subject token position is critical for factual recall. We follow this finding by extracting activations at the last token of each entity name. If ``is a Pok\'{e}mon'' functions partly as a factual association, we would expect it to be concentrated in similar locations---and indeed, we find the strongest signal in middle layers.

\para{Sparse and nonlinear features.}
\citet{cunningham2023sparse} show that sparse autoencoders can decompose model activations into interpretable features, some corresponding to specific entity types. \citet{engels2024not} demonstrate that not all features are one-dimensional: some concepts (e.g., days of the week) are represented on multi-dimensional manifolds. These findings suggest that Pok\'{e}mon-ness could, in principle, be a multi-dimensional feature encoding subtypes (fire, water, etc.) or generations. Our results show that even a single linear direction captures substantial structure, but we acknowledge that richer representations may exist in higher-dimensional subspaces.